% Part: propositional-logic
% Chapter: propositional-logic
% Section: preliminaries

\documentclass[../../../include/open-logic-section]{subfiles}

\begin{document}

\olfileid{pl}{syn}{pre}

\olsection{预备知识}

\begin{thm}[\emph{公式归纳原理}]
  \ollabel{thm:induction}
  若某个性质~$P$ 对所有原子公式成立，且满足条件：
  \begin{enumerate}
    \tagitem{prvNot}{只要它对~$!A$ 成立，则对 $\lnot !A$ 也成立；}{}
    \tagitem{prvAnd}{只要它对 $!A$ 和~$!B$ 成立，则对 $(!A \land !B)$ 也成立；}{}
    \tagitem{prvOr}{只要它对 $!A$ 和~$!B$ 成立，则对 $(!A \lor !B)$ 也成立；}{}
    \tagitem{prvIf}{只要它对 $!A$ 和~$!B$ 成立，则对 $(!A \lif !B)$ 也成立；}{}
    \tagitem{prvIff}{只要它对 $!A$ 和~$!B$ 成立，则对 $(!A \liff !B)$ 也成立；}{}
  \end{enumerate}
  那么 $P$ 对所有公式都成立。
\end{thm}

\begin{proof}
  令 $S$ 为所有具有性质~$P$ 的公式构成的集合。显然 $S \subseteq \Frm[L_0]$。$S$ 满足 \olref[fml]{defn:formulas} 中的所有条件：它包含所有原子公式并且在逻辑算子下封闭。而 $\Frm[L_0]$ 是满足这些条件的最小类，因此 $\Frm[L_0] \subseteq S$。所以 $\Frm[L_0] = S$，即每个公式都具有性质~$P$。
\end{proof}

\begin{prop}\ollabel{prop:balanced}
   $\Frm[L_0]$ 中的任何公式都是\emph{平衡的}，即其左括号与右括号数量相同。
\end{prop}

\begin{prob}
证明 \olref[pl][syn][pre]{prop:balanced}
\end{prob}

\begin{prop} \ollabel{prop:noinit}
公式的任何真前段都不是公式。
\end{prop}

\begin{prob}
  证明 \olref[pl][syn][pre]{prop:noinit}
\end{prob}

\begin{prop}[唯一可读性]
$\Frm[L_0]$ 中的任何公式~$!A$ 都可以唯一地解析为下列形式之一：
\begin{enumerate}
\tagitem{prvFalse}{$\lfalse$。}{}

\tagitem{prvTrue}{$\ltrue$。}{}

\item $\Obj p_n$，其中 $\Obj p_n \in  \PVar$。
  
\tagitem{prvNot}{$\lnot !B$，其中 $!B$ 为某个公式。}{}

\tagitem{prvAnd}{$(!B \land !C)$，其中 $!B$ 和~$!C$ 为某个公式。}{}

\tagitem{prvOr}{$(!B \lor !C)$，其中 $!B$ 和~$!C$ 为某个公式。}{}

\tagitem{prvIf}{$(!B \lif !C)$，其中 $!B$ 和~$!C$ 为某个公式。}{}

\tagitem{prvIff}{$(!B \liff !C)$，其中 $!B$ 和~$!C$ 为某个公式。}{}
\end{enumerate}
此外，这种解析是\emph{唯一}的。
\end{prop}

\begin{proof}
对~$!A$ 进行归纳。例如，假设 $!A$ 有两种不同的解读，即 $(!B \lif !C)$ 和 $(!B' \lif !C')$。那么 $!B$ 和 $!B'$ 必须相同（否则其中一个将是另一个的真前段）；因此，若 $!A$ 的两种解读不同，必然是因为 $!C$ 和 $!C'$ 是同一符号序列的不同解读，而根据归纳假设，这是不可能的。
\end{proof}

\begin{defn}[一致代入]
如果 $!A$ 和 $!B$ 是公式，且 $\Obj p_i$ 是一个命题变元，则 $\Subst{!A}{!B}{\Obj p_i}$ 表示将 $!A$ 中 $\Obj p_i$ 的每一处出现替换为 $!B$ 所得的结果；类似地，将 $\Obj p_1$, \dots,~$\Obj p_n$ 同时替换为公式 $!B_1$, \dots,~$!B_n$ 记作 $\SSubst{!A}{\subst{!B_1}{\Obj p_1},\dots,\subst{!B_n}{\Obj p_n}}$。
\end{defn}

\begin{prob} 对下列五个公式，分别判断其能否表示为代入 \( \Subst{!A}{!B}{\Obj p_i} \)，其中 \( !A \) 为 (i) \( \Obj p_0 \)；(ii) \( ( \lnot \Obj p_0 \land \Obj
 p_1) \)；及 \( ( ( \lnot \Obj p_0 \lif \Obj p_1 ) \land \Obj
 p_2 ) \)。若能，请指明相应的代入。
  \begin{enumerate}
    \item \( \Obj p_1 \)
    \item \( ( \lnot \Obj p_0 \land \Obj p_0 ) \)
    \item \( ( ( \Obj p_0 \lor \Obj p_1 ) \land \Obj p_2 )  \)
    \item \( \lnot ( ( \Obj p_0 \lif \Obj p_1 ) \land \Obj p_2 ) \)
    \item \( (( \lnot ( \Obj p_0 \lif \Obj p_1 ) \lif ( \Obj p_0 \lor \Obj p_1 )) \land \lnot ( \Obj p_0 \land \Obj p_1 )) \)
  \end{enumerate}
\end{prob}

\begin{prob}
用归纳法给出 $\Subst{!A}{!B}{p}$ 的一个数学上严格的定义。
\end{prob}

\end{document}